\documentclass{article}
\usepackage{graphicx, amsmath} % Required for inserting images

\title{Apostol Chapter 1 - Exercise 7}
\author{Afonso Vilhena Francisco}
\date{22 August 2026}

\begin{document}

\maketitle

\textbf{A rational number} $\mathbf{a/b}$ \textbf{with} $\mathbf{(a,b)=1}$ \textbf{is called a} \textbf{\textit{reduced fraction}. if the sum of two reduced fractions is an integer, say} $\mathbf{(a/b)+(c/d)=n},$ \textbf{prove that} $\mathbf{|b|=|d|}$.  

\vspace{0.5cm}

Suppose that $\frac{a}{b}+\frac{c}{d}=n$, with $(a,b)=1=(c,d)$.

Multiplying the previous equation by $bd$, we get

\begin{align}
    ad+cb=bdn
\end{align}

This way we can write both

\begin{align}
    ad=bdn-cb\\
    cb=bdn-ad
\end{align}

Since $(a,b)=1$, we can write $1=ax+by$, for some integers $x,y$. Multiplying it by $d$ and using $(2)$, we get,

\begin{align*}
    d &= adx+bdy\\
      &= (bdn-cb)x+bdy\\
      &= (dnx-cx+dy)b,\\
\end{align*}
that is, $b|d$.

A similar argument shows that $d|b$. Indeed, since $(c,d)=1$, we can write $1=cz+dw$, for some integers $z,w$. Again, multiplying it by $b$ and using $(3)$, it follows that,

\begin{align*}
    b &= cbz+dbw\\
      &= (bdn-ad)z+dbw\\
      &= (bnz-az+bw)d,
\end{align*}
that is, $d|b$.

Finally, since $b|d$ and $d|b$, by item \textit{(j)} of \textit{Theorem 1.1}, follows that $|b|=|d|$
\end{document}
